\par\noindent\textbf{(1)} 均匀背景下动量线化 ρ∂\_t u\_x=(B₀/μ₀)∂\_z b\_x+ρν∂\_z²u\_x，感应方程 ∂\_t b\_x=B₀∂\_z u\_x+η\_m∂\_z²b\_x。磁张力项正比背景场与扰动场的空间导数，两个方程耦合且保证无耗散时能量在两种形式之间交换。
\par\noindent\textbf{(2)} 代入 exp[i(kz−ωt)]，构成矩阵 [−iω+νk²,−ik B₀/(μ₀ρ);−ikB₀,−iω+η\_mk²]。行列式为零得 (ω+iνk²)(ω+iη\_mk²)=k²v\_A²，v\_A=B₀/√(μ₀ρ)。因此 ω\_±=−i(ν+η\_m)k²/2±√[v\_A²k²−(ν−η\_m)²k⁴/4]。
\par\noindent\textbf{(3)} 当 |ν−η\_m|k≪2v\_A，频率实部为±v\_Ak加高阶修正，虚部 −(ν+η\_m)k²/2 给出振幅指数衰减。理想 η\_m=ν=0 时 b\_x=∓√(μ₀ρ)u\_x，磁能 b²/(2μ₀)与动能ρu²/2相等；含耗散时相位关系略偏离。
\par\noindent\textbf{(4)} 若ν≠η\_m，根号非负对应 k≤2v\_A/|ν−η\_m| 的传播分支，超过阈值本征值均纯虚，表示过阻尼；若ν=η\_m，任意k仍有振荡实部但衰减可占主导。k→0两种衰减均趋零，不能在此极限忽略磁耦合后误判完全无波。
